% File: lab-report.tex
\documentclass[11pt]{article}

\usepackage[utf8]{inputenc}
\usepackage[spanish]{babel}
\usepackage[T1]{fontenc}

\usepackage{amssymb} % AMS fonts
\usepackage{amsmath}
\usepackage{amsthm}

\usepackage{../common/file-config}
\usepackage{../common/colours}
\usepackage{../common/listings-config}

\usepackage{graphicx}
% \graphicspath{{res/}}

% Tables
\usepackage{booktabs}
\usepackage[table]{xcolor}
\usepackage{multirow}

% Algorithms
\usepackage{algorithm}
\usepackage{algpseudocode}
\floatname{algorithm}{}
\renewcommand{\thealgorithm}{}

\usepackage{url}
\usepackage[pdftex, colorlinks=true, citecolor=cyan, linkcolor=black, 
urlcolor=black]{hyperref}

\usepackage[backend=biber, style=numeric]{biblatex}
\addbibresource{references.bib}

\begin{document}
    
\begin{titlepage}
	\centering
    % \includegraphics[width=0.5\textwidth]{logo.pdf} \par \vspace{1cm}
	{\scshape \LARGE Asignatura \par}
	\vspace{1cm}
	{\huge\bfseries Práctica... \par}
	\vspace{2cm}
	{\Large Autor} \par
    \texttt{NIP} \par
	\vspace{2cm}
    \textbf{Grupo: }A
    \vfill
    {\Large \today \par}
\end{titlepage}

% \clearemptydoublepage

\thispagestyle{empty}
\tableofcontents

\newpage
% \clearemptydoublepage

\section{Sección 1}

\subsection{\texttt{mdframed}}

\begin{figure}[h!] 
    \center
    \begin{mdframed}[style=code, nobreak=false]
\begin{verbatim}
function Simulated_Annealing(T0, k, Tfinal)
    Sactual <- Solución inicial
    for T=T0 to Tfinal do
        if T=0 then return Sactual
        else Snuevo <- Nueva solución aleatoria
        Inc(C) <- Coste(Snuevo) - Coste(Sactual)
        if (Inc(C) > 0) then Sactual <- Snuevo
        else Sactual <- Snuevo con probabilidad Exp(Inc(C) / F(T))
                                   donde F(T) = k*Exp(-delta*T)
        T <- Evolucion(T)
\end{verbatim}
    \end{mdframed}
    \caption{Pseudocódigo del algoritmo \textit{Simulated Annealing}}
\end{figure}

\subsection{Algoritmos}

\begin{algorithm}
\caption{ECCTA$(m, n, t_1, \dots, t_n, v_1, \dots v_m)$}
    \begin{algorithmic}[1]
        \For{$j \gets 1$ \textbf{to} $m$}
            \State $\hat{W}(j) = 0$; $\hat{A}(j) = \emptyset$
        \EndFor
        
        \For{$i \gets 1$ \textbf{to} $n$}
            \State Sea $j$ el que tiene menor $\hat{W}(j) + t_i/v_j$
            \State Añade $i$ a $\hat{A}(j)$
            \State $\hat{W}(j) = \hat{W}(j) + t_i/v_j$
        \EndFor
        
        \State $\hat{c} = \max_{j} \hat{W}(j)$
        
        \State \Return $\hat{A}$, $\hat{W}$, $\hat{c}$
    \end{algorithmic}
\end{algorithm}

\newpage

\thispagestyle{empty}
\nocite{*}
\printbibliography

\newpage

\appendix

\section{Gestión de la práctica}

\begin{table}[h!]
    \center
    \small
    \begin{tabular}{|p{5cm}|p{5cm}|c|c|c|}
        \bottomrule
        \rowcolor{emerald}
        \textbf{Parte} & \textbf{Tarea} & \textbf{Autor 1} & \textbf{Autor 2} & \textbf{Total}\\
        \toprule
        \bottomrule
        \multirow{5}{*}{\textbf{Análisis de requisitos I}} & Planteamiento del trabajo & 2h & 2h & 4h \\
        \cline{2-5}
                                                           & Recopilación de información & 4h & 5h & 9h \\
        \cline{2-5}
                                                           & Análisis etnográfico & 2h & - & 2h \\
        \cline{2-5}
                                                           & Estudio de la audiencia & 2h & - & 2h \\
        \cline{2-5}                       
                                                           & Estudio de la competencia & 2h & 1h & 3h \\
        \hline
        \multirow{6}{*}{\textbf{Análisis de requisitos II}} & Reunión de control & 2h & 2h & 4h \\
        \cline{2-5}
                                                            & Requisitos funcionales & 1h & 3h & 4h \\
        \cline{2-5}
                                                            & Casos de uso & 2h & 1h & 3h \\
        \cline{2-5}
                                                            & Objetivos de usabilidad & 2h & 1h & 3h \\
        \cline{2-5}                       
                                                            & Estudio de los posibles dispositivos & - & 2h & 2h \\
        \hline
        \multirow{4}{*}{\textbf{Documento de diseño I}} & Reunión de control & 3h & 3h & 6h \\
        \cline{2-5}
                                                        & Primer prototipo & 22h & 8h & 30h \\
        \cline{2-5}
                                                        & Diseño de la Interfaz: Reglas & 3h & 1h & 4h \\
        \cline{2-5}
                                                        & Mapa de navegación & 1h & 4h & 5h \\
        \hline
        \multirow{4}{*}{\textbf{Documento de diseño II}} & Reunión de control & 1h & 1h & 2h \\
        \cline{2-5}
                                                         & Primera evaluación del prototipo & 2h & 4h & 6h \\
        \cline{2-5}
                                                         &  Correcciones al prototipo & 5h & 2h & 7h \\
        \hline
        \multirow{3}{*}{\textbf{Implementación I}} & Reuniones de control & 3h & 3h & 6h \\
        \cline{2-5}
                                                   & Prototipo en Axure & 10h & 20h & 30h \\
        \cline{2-5}
                                                   & Reglas de diseño & 1h & - & 1h \\
        \hline
    
        \multirow{3}{*}{\textbf{Implementación II}} & Reuniones de control & 3h & 3h & 6h \\
        \cline{2-5}
                                                    & Evaluación por heurísticas  & 2h & 2h & 4h \\
        \cline{2-5}
                                                    & Correcciones al prototipo & 1h & 1h & 2h \\
        \toprule
        \bottomrule
        \textbf{Total} & -  & 76h & 69h & 141h \\
        \toprule
    \end{tabular}
    \caption{Horas de trabajo por integrante del equipo y tarea}
\end{table}

\newpage

\section{Código}

\begin{figure}[h!]
    \center
    \begin{lstlisting}[language=Haskell]
    navigate :: Int -> Menu -> Menu
    navigate num menu = menu {pos = newPos, page = head $ newData newPos, firstIndex = (newData newPos) !! 1, lastIndex = last $ newData newPos}
        where
            newPos :: Int
            newPos = if (pos menu) + num < 0 || (menuContent menu) == [] || (pos menu) + num == length (menuContent menu)
                then (pos menu)
                else (pos menu) + num
    
            -- [page, firstIndex, lastIndex]
            newData :: Int -> [Int]
            newData pos = [newPage] ++ newIndexes newPage
                where
                    newPage :: Int
                    newPage
                        | pos < (firstIndex menu) = (page menu) - 1
                        | pos >= (lastIndex menu) = (page menu) + 1
                        | otherwise = (page menu)
    
                    newIndexes :: Int -> [Int]
                    newIndexes page = [page * (numPos menu), (page + 1) * (numPos menu)]
    \end{lstlisting}
    \caption{Código en \texttt{Haskell}}
\end{figure}

\end{document}
